%% ============================================================
%% Material Suplementar
%% Benchmarking multiparadigma para avaliação integrada da
%% qualidade do solo sob manejo conservacionista em Argissolo tropical
%% ============================================================
\documentclass[12pt]{article}
\usepackage[a4paper,margin=2.5cm]{geometry}
\usepackage[utf8]{inputenc}
\usepackage[T1]{fontenc}
\usepackage[brazil]{babel}
\usepackage{booktabs}
\usepackage{multirow}
\usepackage{longtable}
\usepackage{array}
\usepackage{amsmath}
\usepackage{caption}
\usepackage[hidelinks]{hyperref}

\captionsetup{labelfont=bf, font=small}
\renewcommand{\thetable}{S\arabic{table}}

\begin{document}

\begin{center}
{\Large\bfseries Material Suplementar}\\[6pt]
{\large Benchmarking multiparadigma para avaliação integrada da qualidade\\
do solo sob manejo conservacionista em Argissolo tropical}\\[12pt]
\end{center}

\noindent Este documento contém quatro tabelas suplementares referenciadas no manuscrito principal.

%% ============================================================
%%  TABELA S1 �?" Autovalores PCA e pesos dos componentes
%% ============================================================
\clearpage
\begin{table}[ht]
\centering
\caption{Autovalores, variância explicada e pesos dos componentes da PCA para o paradigma ANOVA/MDS (34 indicadores). Apenas os oito componentes retidos (autovalor $> 1$) são apresentados; juntos explicam 83,56\% da variância total.}\label{tab:s1_eigen}
\begin{tabular}{lcccc}
\toprule
Componente & Autovalor & Var.\ expl.\ (\%) & Acumul.\ (\%) & Peso$^a$ \\
\midrule
PC1  & 8,230 & 24,21 & 24,21 & 0,310 \\
PC2  & 5,695 & 16,75 & 40,96 & 0,215 \\
PC3  & 3,284 & 9,66  & 50,61 & 0,124 \\
PC4  & 2,934 & 8,63  & 59,24 & 0,111 \\
PC5  & 2,566 & 7,55  & 66,79 & 0,097 \\
PC6  & 2,435 & 7,16  & 73,95 & 0,092 \\
PC7  & 1,794 & 5,28  & 79,23 & 0,068 \\
PC8  & 1,471 & 4,33  & 83,56 & 0,056 \\
\bottomrule
\end{tabular}
\smallskip

\noindent{\footnotesize $^a$Peso = autovalor$_i$ / $\sum_{i=1}^{8}$autovalor$_i$; utilizado no cálculo aditivo do IQS$_{\text{ANOVA}}$.}
\end{table}

%% ============================================================
%%  TABELA S2 �?" Base completa de regras fuzzy
%% ============================================================
\clearpage
\begin{table}[ht]
\centering
\caption{Base completa de regras do sistema de inferência fuzzy hierárquico Mamdani (IPACSA). Funções de pertinência (FPs): B\,=\,Baixa, M\,=\,Média, A\,=\,Alta; 0\,=\,não importa. Conexão: AND (t-norma mínimo). Defuzzificação: centroide.}\label{tab:s2_fuzzy}
\smallskip
\begin{tabular}{p{4.5cm}lp{5cm}cc}
\toprule
Subíndice (entradas) & Regra & Padrão do antecedente & Saída & Peso \\
\midrule
\multicolumn{5}{l}{\textit{COS -- Conservação do Solo (7 entradas: DMG, DMP, RMP, Ds, Na, ICV, COS$_{\text{estoque}}$)}} \\
 & E1 & M, M, B, B, B, A, A & A & 1,0 \\
 & E2 & M, M, M, M, M, M, M & M & 1,0 \\
 & E3 & B, B, A, A, A, B, B & B & 1,0 \\
 & F1--F21 & Efeito principal$^a$ & B/M/A & 0,1 \\
\midrule
\multicolumn{5}{l}{\textit{ARQ -- Arquitetura de Plantas (6 entradas: ALT, DESP, CESP, NESP, NESPC, PESP)}} \\
 & E1 & A, A, A, A, A, A & A & 1,0 \\
 & E2 & M, M, M, M, M, M & M & 1,0 \\
 & E3 & B, B, B, B, B, B & B & 1,0 \\
 & F1--F18 & Efeito principal$^a$ & B/M/A & 0,1 \\
\midrule
\multicolumn{5}{l}{\textit{STAND -- Estande de Plantas (1 entrada: STAND)}} \\
 & E1 & A & A & 1,0 \\
 & E2 & M & M & 1,0 \\
 & E3 & B & B & 1,0 \\
\midrule
\multicolumn{5}{l}{\textit{IPACSA -- Agregação final (4 entradas: COS, ARQ, STAND, PROD)}} \\
 & E1 & M, A, M, A & A & 1,20 \\
 & E2 & M, M, M, M & M & 1,00 \\
 & E3 & B, B, B, B & B & 0,80 \\
 & E4 & A, M, M, M & A & 1,10 \\
 & E5 & M, A, M, A & A & 1,10 \\
 & E6 & M, M, A, M & A & 1,10 \\
 & E7 & M, A, M, M & A & 1,15 \\
 & F1--F12 & Efeito principal$^a$ & B/M/A & 0,1 \\
\bottomrule
\end{tabular}
\smallskip

\noindent{\footnotesize $^a$Regras de efeito principal: para cada variável de entrada $k$, três regras são geradas onde apenas a $k$-ésima entrada é definida como B, M ou A (demais = 0, i.e., não importa), mapeando diretamente para o nível de saída correspondente. Essas regras (peso\,=\,0,1) previnem saídas NaN em regiões do espaço de entrada não cobertas pelas regras de especialista.}
\end{table}

%% ============================================================
%%  TABELA S3 �?" Diagnósticos PLS-SEM
%% ============================================================
\clearpage

\begin{table}[ht]
\centering
\caption{Modelo estrutural PLS-SEM: coeficientes de caminho, intervalos de confiança por bootstrap e ajuste do modelo (cinco construtos, 25 variáveis manifestas após triagem VIF).}\label{tab:s3_paths}
\begin{tabular}{lcccccc}
\toprule
Caminho & $\gamma$ & Média Boot. & DP Boot. & $t$ & IC 95\% & $p$ \\
\midrule
FÍSICO $\to$ REND   & $-$0,122 & $-$0,013 & 0,160 & 0,76  & [$-$0,28;\;0,25]  & 0,992 \\
QUÍMICO $\to$ REND  &    0,097 &    0,104 & 0,102 & 0,95  & [$-$0,14;\;0,28]  & 0,264 \\
MICRO $\to$ REND    &    0,175 &    0,119 & 0,127 & 1,39  & [$-$0,22;\;0,28]  & 0,284 \\
AGRO $\to$ REND     &    0,622 &    0,608 & 0,051 & 12,30 & [\phantom{$-$}0,50;\;0,70] & $<$0,001 \\
\midrule
\multicolumn{7}{l}{$R^2 = 0,690$; $R^2$ Adj.\,=\,0,678; Bootstrap: 500 reamostragens, semente = 42.} \\
\bottomrule
\end{tabular}
\end{table}

\begin{table}[ht]
\centering
\caption{Pesos externos PLS-SEM (Modo~B, formativo) e significância por bootstrap para 25 variáveis manifestas retidas.}\label{tab:s3_weights}
\begin{tabular}{llrrrrl}
\toprule
Construto & Indicador & Peso & DP Boot. & $t$ & IC 95\% & $p$ \\
\midrule
\multirow{9}{*}{FÍSICO}
 & DMG   &  0,068 & 0,192 & 0,35 & [$-$0,29;\;0,41] & 0,792 \\
 & COS   &  0,086 & 0,217 & 0,40 & [$-$0,39;\;0,47] & 0,784 \\
 & Na    &  0,153 & 0,216 & 0,71 & [$-$0,36;\;0,43] & 0,864 \\
 & ICV   &  0,031 & 0,196 & 0,16 & [$-$0,34;\;0,44] & 0,888 \\
 & MI    & $-$0,301 & 0,315 & 0,95 & [$-$0,50;\;0,60] & 0,800 \\
 & MA    &  0,684 & 0,691 & 0,99 & [$-$1,09;\;1,17] & 0,924 \\
 & VIB   &  0,178 & 0,268 & 0,66 & [$-$0,48;\;0,55] & 0,900 \\
 & RP    &  0,314 & 0,374 & 0,84 & [$-$0,66;\;0,68] & 0,980 \\
 & AD/PT &  1,200 & 1,108 & 1,08 & [$-$1,41;\;1,47] & 0,924 \\
\midrule
\multirow{6}{*}{QUÍMICO}
 & pH    & $-$0,365 & 0,206 & 1,77 & [$-$0,61;\;0,42] & 0,092 \\
 & P     & $-$0,508 & 0,263 & 1,93 & [$-$0,75;\;0,52] & 0,088 \\
 & K     & $-$0,044 & 0,242 & 0,18 & [$-$0,51;\;0,44] & 0,936 \\
 & Mg    &  0,096 & 0,196 & 0,49 & [$-$0,32;\;0,44] & 0,640 \\
 & CTC   & $-$0,197 & 0,210 & 0,94 & [$-$0,51;\;0,46] & 0,316 \\
 & EstN  &  0,618 & 0,330 & 1,87 & [$-$0,51;\;1,06] & 0,088 \\
\midrule
\multirow{4}{*}{MICRO}
 & qMIC  & $-$0,580 & 0,445 & 1,30 & [$-$0,98;\;0,80] & 0,276 \\
 & qCO$_2$ &  0,357 & 0,295 & 1,21 & [$-$0,43;\;0,78] & 0,276 \\
 & CCO$_2$ & $-$0,246 & 0,363 & 0,68 & [$-$0,72;\;0,75] & 0,656 \\
 & NMIC  &  0,927 & 0,596 & 1,56 & [$-$0,90;\;1,05] & 0,276 \\
\midrule
\multirow{4}{*}{AGRO}
 & ALT   &  0,601 & 0,089 & 6,74 & [\phantom{$-$}0,40;\;0,76] & $<$0,001 \\
 & DESP  &  0,124 & 0,087 & 1,43 & [$-$0,06;\;0,30] & 0,172 \\
 & CESP  &  0,337 & 0,081 & 4,16 & [\phantom{$-$}0,17;\;0,48] & $<$0,001 \\
 & STAND &  0,529 & 0,108 & 4,90 & [\phantom{$-$}0,32;\;0,72] & 0,004 \\
\midrule
\multirow{2}{*}{REND}
 & NESP  &  0,534 & 0,051 & 10,46 & [\phantom{$-$}0,44;\;0,62] & $<$0,001 \\
 & PESP  &  0,643 & 0,052 & 12,34 & [\phantom{$-$}0,56;\;0,76] & $<$0,001 \\
\bottomrule
\end{tabular}
\smallskip

\noindent{\footnotesize A triagem VIF (limiar = 5) excluiu Ca (VIF\,=\,26,5), MOS (VIF\,=\,9,8), CBM (VIF\,=\,23,9) e PROD (VIF\,=\,20,1). Construto REND estimado com Modo~A (reflexivo); demais com Modo~B (formativo). Bootstrap: 500 reamostragens.}
\end{table}

%% ============================================================
%%  TABELA S4 �?" Importância de variáveis RF
%% ============================================================
\clearpage
\begin{table}[ht]
\centering
\caption{Importância de variáveis do Random Forest (modelo com 34 indicadores, 500 árvores, $m_{\text{try}} = 11$). Variáveis ordenadas por importância de permutação (\%IncMSE).}\label{tab:s4_rf}
\begin{tabular}{rlccc}
\toprule
Posição & Indicador & Domínio & \%IncMSE & IncNodePurity \\
\midrule
1  & NESP  & Rendimento      & 20,22 & 30,46 \\
2  & STAND & Agronômico      & 19,49 & 23,51 \\
3  & NESPC & Rendimento      & 13,19 & 8,44 \\
4  & PROD  & Rendimento      & 13,15 & 8,02 \\
5  & PESP  & Rendimento      & 12,51 & 8,26 \\
6  & DESP  & Agronômico      & 12,33 & 9,38 \\
7  & ICV   & Físico          & 11,19 & 5,95 \\
8  & CESP  & Agronômico      & 11,13 & 5,27 \\
9  & P     & Químico         & 8,04 & 3,01 \\
10 & Ca    & Químico         & 7,63 & 1,32 \\
11 & MOS   & Químico         & 7,35 & 1,68 \\
12 & CTC   & Químico         & 6,68 & 1,48 \\
13 & K     & Químico         & 6,64 & 1,81 \\
14 & MI    & Físico          & 6,47 & 1,16 \\
15 & pH    & Químico         & 6,38 & 1,35 \\
16 & EstN  & Químico         & 6,11 & 1,36 \\
17 & Al    & Químico         & 6,09 & 1,23 \\
18 & qMIC  & Microbiológico  & 5,89 & 1,11 \\
19 & CBM   & Microbiológico  & 5,64 & 0,97 \\
20 & AD/PT & Físico          & 5,51 & 1,20 \\
21 & VIB   & Físico          & 5,39 & 1,16 \\
22 & RB    & Microbiológico  & 5,23 & 1,12 \\
23 & NBM   & Microbiológico  & 5,09 & 1,15 \\
24 & RP    & Físico          & 4,96 & 0,86 \\
25 & Mg    & Químico         & 4,68 & 0,60 \\
26 & RMP   & Físico          & 4,53 & 1,29 \\
27 & qCO$_2$  & Microbiológico & 4,38 & 0,80 \\
28 & MA    & Físico          & 4,31 & 1,00 \\
29 & COS   & Físico          & 4,18 & 1,78 \\
30 & DMG   & Físico          & 3,87 & 1,10 \\
31 & ALT   & Agronômico      & 3,40 & 1,39 \\
32 & DMP   & Físico          & 3,30 & 1,56 \\
33 & Na    & Físico          & 2,97 & 1,10 \\
34 & Ds    & Físico          & 2,95 & 1,05 \\
\bottomrule
\end{tabular}
\smallskip

\noindent{\footnotesize \%IncMSE, aumento percentual no erro quadrático médio por permutação da variável; IncNodePurity, redução total na impureza dos nós (soma de quadrados dos resíduos).}
\end{table}

\end{document}
